\honest{What Is Evidence?}{Eliezer Yudkowsky}{2007}

I open with two quotations about truth, from Tarski and from Aristotle. This post is about evidence. I return to Tarski's snow in my last two paragraphs.

Your shoelaces come untied. Light bounces off them into your eye, your brain processes it, and you come to believe that they are untied. Your belief now mirrors your shoelaces.

Evidence is ``an event entangled, by links of cause and effect, with whatever you want to know about.'' A machine that beeps for every lottery number, winning or losing, tells you nothing. Evidence must come out differently depending on the facts. A parenthesis gives the technical version: mutual information, relative to your current uncertainty. \nb{The first definition speaks of causes in the world; the second is measured relative to what you already believe, so the same event can be evidence for one person and not for another.}

Light from your shoelaces leaves no trace in a rock it hits, which is why rocks make poor witnesses in court. A photograph or a working eye keeps a trace. If your eyes and brain work, you become tied to your shoelaces.

Next, rationalists value the ``paradoxical-seeming'' claim that a belief is worth holding only if something could change it. If your retina ended up in the same state whatever light entered it, you would be blind. Some belief systems, ``in a rather obvious trick,'' say that certain beliefs must be held no matter what you see. Hence ``blind faith.'' A belief that does not depend on what you see blinds you as surely as losing your eyes. I do not name the belief systems or quote any of them.

``Rational thought produces beliefs which are themselves evidence.''

Your beliefs, once spoken, can be evidence for someone else, since speaking and hearing are causes too. Saying ``My shoelaces are untied'' over the phone shares your tie to your shoelaces with a friend. So rational beliefs are ``contagious, among honest folk who believe each other to be honest.'' Therefore a belief held for private reasons that cannot be passed on is ``suspicious.''

If my own model says my beliefs should not be contagious to others, I should ``stop believing.'' And if you understand this ``on a gut level,'' you will stop ``automatically,'' because a belief that is not tied to reality is not accurate. \nb{A stopped clock is not tied to the time, and it is right twice a day. Unreliable is not the same as false.}

So try to explain why your way of thinking produces beliefs that mirror reality: why minds that think as you do will believe ``snow is white'' if and only if snow is white. If you do not believe your thinking is tied to reality, why believe what it tells you? \nb{The explanation would have to come from that same way of thinking. The post does not say so.}
